
\textbf{Théorème 1.13.}
\emph{Soient @@M0@@ un schéma noethérien et @@M1@@ un @@M2@@-schéma projectif et muni d'un faisceau ample @@M3@@ dont les induits @@M4@@, avec @@M5@@, sont engendrés par leur @@M6@@; soit @@M7@@ une famille de classes de faisceaux cohérents sur les fibres de @@M8@@. Alors les conditions suivantes sont équivalentes} :
\begin{enumerate}[label=(\roman*)]
\item \emph{@@M9@@ est limitée. De plus, si tous les @@M10@@ sont localement libres de rang @@M11@@, alors on peut supposer que @@M12@@ est limitée par un faisceau @@M13@@ sur un @@M14@@, avec @@M15@@ localement libre de rang @@M16@@.}
\item \emph{L'ensemble des polynômes de Hilbert @@M17@@ des @@M18@@ est fini, et il existe une suite d'entiers @@M19@@ telle que @@M20@@ soit une @@M21@@-famille.}
\item \emph{L'ensemble des polynômes de Hilbert @@M22@@ des @@M23@@ est fini, et il existe un entier @@M24@@ tel que tous les @@M25@@ soient @@M26@@-réguliers.}
\item \emph{L'ensemble des polynômes de Hilbert @@M27@@ des @@M28@@ est fini, et @@M29@@ est contenue dans la famille des classes de quotients de faisceaux de la forme @@M30@@, où @@M31@@ est un faisceau cohérent sur un @@M32@@. De plus, on peut supposer @@M33@@ et @@M34@@ de la forme @@M35@@.}
\item \emph{@@M36@@ est contenue dans la famille des classes de conoyaux d'homomorphismes de la forme @@M37@@ où @@M38@@ et @@M39@@ sont des faisceaux cohérents sur un @@M40@@. De plus, on peut supposer @@M41@@ et @@M42@@ et @@M43@@ de la forme @@M44@@ et @@M45@@.}
\end{enumerate}
En effet, supposons que @@M46@@ soit limitée par un faisceau @@M47@@ sur un @@M48@@; quitte à couper @@M49@@ en morceaux, en utilisant le théorème de platitude générique (EGA IV 6.9.1), on peut supposer @@M50@@ plat sur @@M51@@. Alors,
